\begin{table}[ht]
\centering
\caption{Recovery of residential population for community water systems one HUC12 directly downstream of coal mining.}
\label{tab:backcast_recovery}
\begin{adjustbox}{max width=\textwidth}
\begin{tabular}{lrl}
\toprule
& Systems & Notes \\
\midrule
\multicolumn{3}{l}{\textit{Panel A: Coverage of the downstream roster}} \\
Ever downstream of a mine, 1983--2024 & 367 & Target set, $\mathcal{D}=\bigcup_t \mathcal{D}_t$ \\
\quad Matched to an EPA SABS service-area polygon & 239 & Preferred tier \\
\quad Matched to a county polygon (fallback) & 127 & 57 distinct counties \\
\quad Unmatched (no polygon, no county) & 1 & Dropped \\
\textbf{Population recovered} & \textbf{366} & 99.7\% of the roster, across 18 states \\
\midrule
\multicolumn{3}{l}{\textit{Panel B: Years recovered}} \\
Backcast window & 1990--2024 & 35 annual observations per system \\
\quad Decennial anchor years & 4 & 1990, 2000, 2010, 2020 (block apportionment) \\
\quad Intercensal years & 31 & PEP county growth, re-anchored each decade \\
System $\times$ year panel & 12,810 & Balanced, no missing values \\
Roster years not recovered & --- & 1983--1989: no nationwide block geography \\
\bottomrule
\end{tabular}
\end{adjustbox}
\begin{minipage}{\textwidth}\footnotesize
\vspace{0.5em}
\textit{Notes:} A community water system is downstream if at least one of its intakes lies in a HUC12 directly downstream of a HUC12 with an active coal mine, and none lies in a mining HUC12 itself. Residential population inside each system's exposure polygon is apportioned from decennial Census blocks by areal weights and carried to intercensal years with county Population Estimates Program growth. The 1990 floor binds because nationwide digital block boundaries begin with the 1990 TIGER/Line release.
\end{minipage}
\end{table}
